% Template:     Reporte LaTeX
% Documento:    Replication of Barlas Fig 11.15
% Codificación: UTF-8

\documentclass[
	english,
	letterpaper, oneside
]{article}

% INFORMACIÓN DEL DOCUMENTO
\def\documenttitle {Replicating Figure 11.15: A Thread-Configuration Sweep of the Hybrid CPU+GPU Mandelbrot Renderer}
\def\documentsubtitle {}
\def\documentsubject {Comparing the Barlas (2022) load-balancing case study on an i7-9750H + GTX 1660 Ti Max-Q laptop}
\def\documentdate {\today}

\def\documentauthor {Matias Saavedra}
\def\coursename {Multicore and GPU Programming}
\def\coursecode {CISC 701}

\def\universityname {Universidad de Chile}
\def\universityfaculty {Facultad de Ciencias Físicas y Matemáticas}
\def\universitydepartment {Departamento de Ciencias de la Computación}
\def\universitydepartmentimage {departamentos/dcc}
\def\universitylocation {Santiago de Chile}

../profiling_report/template.tex

\begin{document}

\templatePagecfg
\templateFinalcfg

\inserttitle

% -----------------------------------------------------------------------
\begin{abstract}
This report documents a follow-up profiling experiment on
\texttt{mandelHybrid}, the hybrid CPU+GPU Mandelbrot renderer from
Barlas, \emph{Multicore and GPU Programming} (2nd ed., 2022),
Section~11.5.1. The goal was to reproduce Figure~11.15 of that book---a
sweep of end-to-end execution time across seven thread/GPU
configurations---on a different machine (an i7-9750H~+~GTX~1660~Ti~Max-Q
laptop) and compare the resulting curve to the book's Ryzen
3900X~+~RTX~2060~Super result. A small instrumentation patch was added
to the existing harness: an end-to-end \texttt{clock\_gettime} wall-clock
timer, a self-describing configuration banner, and two boolean flags
(\texttt{profileQuiet}, \texttt{profileSave}) that suppress per-region
\texttt{stderr} output and skip PNG writes when batch-running. A Bash
driver (\texttt{sweep.sh}) executes seven configurations three times each
into a CSV; a Python script (\texttt{plot\_fig1115.py}) renders a chart
in the style of Figure~11.15. The replication shows the same monotonic
\emph{decrease} as one moves from GPU-alone through hybrid configurations
with more CPU workers, but the slope is much shallower than in the book:
on this laptop the GPU is 2.1$\times$ faster than the CPU at full
parallelism, so adding 11~CPU workers buys only a 30\% speed-up over
GPU-alone, versus the $\sim$50\% the book observed when its CPU and GPU
were roughly on par.
\end{abstract}

% -----------------------------------------------------------------------
\section{Background and Goals}

Figure~11.15 of Barlas (2022) plots the average end-to-end execution time
of a 100-frame, full-HD Mandelbrot animation as a function of thread
configuration. Eight bars are shown, corresponding to: CPU-only with the
full thread count of the machine, GPU-only, and six hybrid
configurations with $1, 2, 4, 8, 16, 23$ pure CPU workers running
alongside one GPU thread. A line overlay reports each bar's time as a
percentage of the GPU-alone time, anchoring the comparison.

The book's measurements were taken on an AMD Ryzen 3900X (12 cores,
24~threads) paired with an NVIDIA RTX 2060 Super. On that platform the
CPU and GPU are roughly on par for the Mandelbrot iteration kernel and
the hybrid run cuts execution time approximately in half relative to
either alone. This report addresses three questions:

\begin{enumerate}
    \item What instrumentation is needed to record the metric the figure
    actually plots, given the existing per-region NVTX/CUDA-event
    profiling captures only \emph{components} of total time and not the
    end-to-end wall clock?
    \item How does the same workload behave on an Intel Core i7-9750H
    (6~cores, 12~threads) paired with an NVIDIA GeForce GTX~1660~Ti
    Max-Q---a laptop GPU that is roughly comparable in fp64 throughput
    to the RTX 2060 Super but mated to a CPU that is half as wide?
    \item What is the shape of the curve and what does any deviation
    from the book's curve reveal about the CPU/GPU balance of this
    machine?
\end{enumerate}

\subsection{What Was Already Instrumented}

A prior profiling pass had already added NVTX annotations and CUDA
event timing to the program (documented in
\texttt{profiling\_report/main.tex}). Those mechanisms give:

\begin{itemize}
    \item per-region GPU kernel time and device-to-host \texttt{memcpy}
    time, summed over the run and reported in
    \texttt{CUDAmemCleanup()};
    \item per-region CPU compute time accumulated by a
    \texttt{thread\_local} pair of counters in
    \texttt{mandelregion.cpp} and reported in \texttt{printCPUSummary()};
    \item nested NVTX ranges (\texttt{examine region},
    \texttt{CPU region compute}, \texttt{GPU region compute}) consumable
    by Nsight Systems.
\end{itemize}

What was \emph{not} captured was the total end-to-end wall-clock time of
the program---the makespan---which is exactly the quantity plotted in
Figure~11.15. \texttt{cpuTotalMs} excludes idle time. \texttt{totalKernelMs}
covers only the GPU's busy stretches. The thread summaries are
independent and the final time is determined by whichever thread
finishes last, plus any image-save tail. Two trivial routes existed:
wrap the binary with \texttt{/usr/bin/time -f "\%e"}, or bake a
\texttt{clock\_gettime} pair into \texttt{main.cpp}. The latter was
chosen because it makes the binary self-reporting (easier for batch
scripts to grep) and because it gives a clean point to also emit a
configuration banner.

% -----------------------------------------------------------------------
\section{Instrumentation Changes}

\subsection{Two Global Flags Defined in \texttt{main.cpp}}

Two booleans were added with \texttt{extern "C"} linkage so that the
\texttt{nvcc}-compiled \texttt{kernel.o} and the \texttt{g++}-compiled
\texttt{.o} files can share the symbol without name-mangling concerns:

\begin{sourcecode}{cpp}{main.cpp --- profiling flag definitions.}
// Profiling flags consulted by kernel.cu, mandelregion.cpp and
// mandelframe.cpp.  Defined here so a single command-line switch
// controls them.  C linkage keeps the symbol name identical in nvcc
// and g++ object files.
extern "C" int profileQuiet = 0;   // 1 = suppress per-region stderr prints
extern "C" int profileSave  = 1;   // 0 = skip PNG saves (pure-compute timing)
\end{sourcecode}

Two new optional positional command-line arguments were added so the
flags can be flipped without recompiling:

\begin{sourcecode}{cpp}{main.cpp --- argument parsing.}
// Usage: spec_file [numThr] [GPUenable] [diffT] [pixT] [quiet] [save]
if (argc > 6) profileQuiet = atoi(argv[6]);
if (argc > 7) profileSave  = atoi(argv[7]);
\end{sourcecode}

\subsection{Per-Region Print Gates}

Three call sites had previously emitted one \texttt{stderr} line per
leaf region. Across the 4{,}603 leaves of the standard run, plus the
banners, an unmuted execution emits roughly 5{,}000 lines per run; 7
configurations $\times$ 3 reps would produce $\sim$100{,}000 lines of
stderr noise that dwarfs the summary lines we actually want to parse.
Each call site was wrapped in an \texttt{if (!profileQuiet)} check, with
the per-thread \texttt{printCPUSummary()} and the
\texttt{CUDAmemCleanup()} summary deliberately left always on:

\begin{sourcecode}{cuda}{kernel.cu --- GPU per-region print gated.}
// Defined in main.cpp.  Suppresses per-region stderr lines when set,
// leaving the once-per-run summary intact.
extern "C" int profileQuiet;

// ... inside hostFE() ...
if (!profileQuiet)
  fprintf(stderr,
          "[GPU region %4ld] size %4dx%4d  kernel %.3f ms  D->H %.3f ms\n",
          totalRegions, resX, resY, kernelMs, memcpyMs);
\end{sourcecode}

\begin{sourcecode}{cpp}{mandelregion.cpp --- CPU per-region print gated.}
extern "C" int profileQuiet;

// ... inside compute(), CPU path ...
if (!profileQuiet)
  fprintf(stderr,
          "[CPU region %4ld] size %4dx%4d  compute %.3f ms\n",
          cpuRegionCount, pixelsX, pixelsY, ms);
\end{sourcecode}

\subsection{Skip-Save Gate}

Figure~11.15 in the book includes the PNG-encoding tail; for
side-by-side comparison the sweep also defaults to \texttt{save=1}, but
the option to time pure compute is useful for follow-up analysis. The
guard was added inside \texttt{MandelFrame::regionComplete}, the single
point at which the atomically-decremented remaining-regions counter
hits zero and the frame is written:

\begin{sourcecode}{cpp}{mandelframe.cpp --- PNG-save gate.}
// Defined in main.cpp.  When 0, regionComplete skips img->save() so we
// can time the pure compute path without the PNG-encode tail.
extern "C" int profileSave;

void MandelFrame::regionComplete()
{
  if (remainingRegions.fetchAndAddOrdered(-1) == 1)  // last became 0
  {
     if (profileSave)
       img->save(fname, "PNG");
  }
}
\end{sourcecode}

\subsection{End-to-End Wall-Clock Timer and Configuration Banner}

The new timer brackets exactly the compute phase: it starts immediately
before \texttt{CUDAmemSetup} (when GPU is enabled) and stops after the
\texttt{wait()} on every worker thread returns. Image saves are
naturally included because they occur on the worker thread that
decrements \texttt{remainingRegions} to 0, before that worker exits its
\texttt{run()} loop and is joined.

\begin{sourcecode}{cpp}{main.cpp --- banner, timer, machine-parseable output.}
fprintf(stderr,
        "[config] numThr=%d gpuEnable=%d diffT=%.3f pixT=%d "
        "res=%dx%d frames=%d quiet=%d save=%d\n",
        numThreads, (int)enableGPU, diffT, pixT,
        resolutionX, resolutionY, numframes,
        profileQuiet, profileSave);

// ... thread launch ...

struct timespec t0, t1;
clock_gettime(CLOCK_MONOTONIC, &t0);

if (enableGPU) {
    CUDAmemSetup(resolutionX, resolutionY);
    thr[0]->run();
    CUDAmemCleanup();
} else {
    thr[0]->run();
}
for (int i = 1; i < numThreads; i++)
    thr[i]->wait();

clock_gettime(CLOCK_MONOTONIC, &t1);
double elapsed_s = (t1.tv_sec - t0.tv_sec)
                 + (t1.tv_nsec - t0.tv_nsec) * 1e-9;

// Machine-parseable line: sweep.sh greps for this exact prefix.
fprintf(stderr, "[total_elapsed_s] %.6f\n", elapsed_s);
\end{sourcecode}

\texttt{CLOCK\_MONOTONIC} is used for the same reason as in the
per-region instrumentation: it is immune to wall-clock jumps from NTP
or DST. The marker line \texttt{[total\_elapsed\_s]} is chosen because
it is unique across all stderr output and is therefore safe to extract
with a single \texttt{grep ... \textbackslash{}| awk '\{print \$2\}'}.

\subsection{Build-System Adjustment (nvcc + gcc 16)}

A separate environmental issue surfaced during the rebuild: the host
\texttt{g++} on the test machine had been upgraded to 16.1.1, but
\texttt{nvcc} 13.2 only supports up to GCC~15. The compile failed with
\texttt{type\_traits} errors deep inside libstdc++. The Makefile was
amended to pin the host compiler explicitly via \texttt{-ccbin}:

\begin{sourcecode}{make}{Makefile --- pin nvcc's host compiler.}
# nvcc 13.2 does not support gcc 16 yet; pin to g++-15 for the host
# compiler and for nvcc's -ccbin so both halves of the program share
# an ABI.
HOSTCC    = g++-15
NVCC      = nvcc -ccbin $(HOSTCC)
CC        = $(HOSTCC)
\end{sourcecode}

Using the same compiler for both halves of the build is important
because the C++ runtime types (e.g.\ \texttt{std::atomic\_int},
\texttt{std::vector}) flow across the
\texttt{nvcc}/\texttt{g++} boundary via the shared
\texttt{MandelRegion}/\texttt{MandelFrame} headers. A divergent ABI here
would manifest as silent corruption rather than a link error.

% -----------------------------------------------------------------------
\section{Batch Driver and Plot Script}

\subsection{\texttt{sweep.sh}}

The driver script runs seven configurations $N$ times each and emits a
single CSV row per run. Configurations are tuned for a 12-thread
machine: total threads (the \texttt{numThreads} argument to
\texttt{mandelHybrid}) never exceed 12, and the count of \emph{pure CPU}
workers in a hybrid run is \texttt{numThreads - 1} because thread~0 is
the GPU driver.

\begin{sourcecode}{bash}{sweep.sh --- configuration matrix and main loop.}
CONFIGS=(
  "CPU12         12  0"
  "GPU            1  1"
  "GPU+1CPU       2  1"
  "GPU+2CPU       3  1"
  "GPU+4CPU       5  1"
  "GPU+8CPU       9  1"
  "GPU+11CPU     12  1"
)

CSV="$OUT_ABS/results.csv"
echo "label,numThreads,gpuEnable,rep,elapsed_s" > "$CSV"

for entry in "${CONFIGS[@]}"; do
  set -- $entry
  label=$1; nthr=$2; gpu=$3
  for r in $(seq 1 "$REPS"); do
    rm -rf "$SCRATCH" && mkdir -p "$SCRATCH"
    log="$OUT_ABS/logs/${label}.r${r}.log"
    pushd "$SCRATCH" >/dev/null
    # quiet=1, save passed through (default 1 to match Fig 11.15).
    "$BIN_ABS" "$SPEC_ABS" "$nthr" "$gpu" "$DIFFT" "$PIXT" 1 "$SAVE" \
        > "$log.stdout" 2> "$log.stderr"
    popd >/dev/null
    elapsed=$(grep '^\[total_elapsed_s\]' "$log.stderr" | awk '{print $2}')
    echo "$label,$nthr,$gpu,$r,$elapsed" >> "$CSV"
  done
done
\end{sourcecode}

Two design choices are worth highlighting:

\begin{itemize}
    \item Each run is executed with its own scratch \texttt{cwd}, so
    the PNG files for run~$r$ of config~$c$ do not pollute the next
    invocation and disk usage stays bounded (50\,MB peak for a single
    full-HD 100-frame run, cleaned between runs).
    \item Configurations, reps count, save flag, and thresholds are all
    environment-variable overrides
    (\texttt{REPS}, \texttt{SAVE}, \texttt{DIFFT}, \texttt{PIXT},
    \texttt{SPEC}, \texttt{OUT}) so the same script supports the main
    Fig~11.15 sweep, a no-save (pure-compute) sweep, and a follow-up
    threshold sensitivity study without code changes.
\end{itemize}

\subsection{\texttt{plot\_fig1115.py}}

The plot script reads the CSV, groups by configuration, computes the
mean and population standard deviation of \texttt{elapsed\_s} per group,
and emits a Fig~11.15-style chart: blue bars with error caps on a
left-axis time scale, a red percentage-of-GPU-alone line on the right
axis, numeric labels on every element. Bar value labels are placed
\emph{inside} each bar in white to avoid colliding visually with the
percent labels above the line markers.

\begin{sourcecode}{python}{plot\_fig1115.py --- core rendering block.}
bars = ax1.bar(x, means, yerr=stds, capsize=4,
               color="#4472C4", edgecolor="#1f3a73",
               label="Mean execution time (s)")
ax1.set_xticks(x); ax1.set_xticklabels(order, rotation=20, ha="right")
ax1.set_ylabel("Execution time (s)")
ax1.set_ylim(0, max(means) * 1.28)
for rect, m in zip(bars, means):
    ax1.text(rect.get_x() + rect.get_width() / 2,
             rect.get_height() * 0.5,
             f"{m:.2f} s",
             ha="center", va="center", fontsize=9,
             color="white", fontweight="bold")

ax2 = ax1.twinx()
ax2.plot(x, pct, "o-", color="#C00000",
         label=f"% of {baseline_label}-alone")
ax2.set_ylim(0, max(pct) * 1.28)
\end{sourcecode}

% -----------------------------------------------------------------------
\section{Experimental Setup}

\begin{table}[H]
    \centering
    \caption{Test platform compared with the book's.}
    \label{tab:hw}
    \begin{tabular}{lll}
        \hline
        & \textbf{This work} & \textbf{Barlas (2022)} \\
        \hline
        CPU            & Intel Core i7-9750H        & AMD Ryzen 3900X \\
        Cores/threads  & 6 / 12                     & 12 / 24 \\
        GPU            & GTX 1660 Ti Max-Q          & RTX 2060 Super \\
        GPU memory     & 6\,GiB                     & 8\,GiB \\
        Compute cap.   & 7.5 (Turing)               & 7.5 (Turing) \\
        nvcc version   & 13.2                       & not stated \\
        Reps per config & 3                         & 10 \\
        \hline
    \end{tabular}
\end{table}

The workload was identical: 100 frames at $1920 \times 1080$ pixels,
starting at corners $(-1.9, 0.94)$--$(1.2, -0.94)$ with
\texttt{maxIterations}\,=\,100, ending at
$(0.00164, -0.82247)$--$(0.00164, -0.82247)$ with
\texttt{maxIterations}\,=\,10{,}000. The default adaptive-subdivision
parameters were used: \texttt{diffThreshold}\,=\,0.5,
\texttt{pixelSizeThresh}\,=\,32{,}768. PNG saves were enabled
(\texttt{save}\,=\,1), matching the book.

% -----------------------------------------------------------------------
\section{Results}

The CSV emitted by \texttt{sweep.sh} contained 21 rows (7 configs $\times$
3 reps). Summary statistics are in Table~\ref{tab:results}; the chart in
Figure~\ref{fig:replicated} is the direct analogue of the book's
Fig~11.15.

\begin{table}[H]
    \centering
    \caption{End-to-end wall-clock time per configuration, 3 reps each.}
    \label{tab:results}
    \begin{tabular}{lrrrrr}
        \hline
        \textbf{Config} & \textbf{Mean (s)} & \textbf{Std (s)} & \textbf{Min (s)} & \textbf{Max (s)} & \textbf{\% of GPU-alone} \\
        \hline
        CPU12      & 163.470 & 0.392 & 163.018 & 163.702 & 213.6\% \\
        GPU        &  76.543 & 0.166 &  76.360 &  76.683 & 100.0\% \\
        GPU+1CPU   &  68.453 & 0.586 &  68.061 &  69.127 &  89.4\% \\
        GPU+2CPU   &  64.398 & 0.412 &  63.959 &  64.777 &  84.1\% \\
        GPU+4CPU   &  60.427 & 0.692 &  59.697 &  61.073 &  78.9\% \\
        GPU+8CPU   &  56.194 & 0.222 &  56.048 &  56.450 &  73.4\% \\
        GPU+11CPU  &  53.618 & 0.118 &  53.500 &  53.737 &  70.0\% \\
        \hline
    \end{tabular}
\end{table}

\insertimage{img/fig1115_replicated.png}{width=\linewidth}
            {Replicated Fig~11.15 on the i7-9750H + GTX 1660 Ti Max-Q
             laptop.  Blue bars: mean wall-clock execution time over 3
             reps; black caps would denote standard deviation but are
             invisible at this scale because variance is below 1\% of
             the mean for every configuration.  Red line: percentage of
             GPU-alone execution time.  \label{fig:replicated}}

For reference, the book's Fig~11.15 is reproduced in
Figure~\ref{fig:book}. Note the qualitatively similar shape---a
monotonic decrease from CPU-only on the left through GPU-only and into
the hybrid configurations---but the very different ratio between the
left two bars.

\insertimage{img/fig1115_book.png}{width=\linewidth}
            {Reproduction of Fig~11.15 from Barlas (2022).  Note that
             ``CPU'' here is a 24-thread Ryzen 3900X and ``GPU'' is an
             RTX~2060 Super; the percentages on the line are relative
             to the leftmost (24-thread CPU) bar in this rendition.
             \label{fig:book}}

% -----------------------------------------------------------------------
\section{Analysis}

\subsection{Variance Is Negligible}

Per-configuration standard deviations are at most 0.69~s on a 60-second
run (1.1\%); the smallest hybrid configurations land at 0.12--0.22~s
(0.2--0.4\%). Three reps was therefore comfortably enough to anchor
each bar---the error caps in Figure~\ref{fig:replicated} are not even
visible. The 10 reps used in the book offer diminishing returns when
the run-to-run noise is already below 1\%; reps are better spent on
varying parameters than on shrinking already-tight error bars.

\subsection{GPU Dominates on This Machine}

The headline number is CPU12~/~GPU $\approx$ 2.14. The 6-core/12-thread
i7-9750H is 2.1$\times$ slower than the GTX~1660~Ti Max-Q at this
workload, despite running at full thread parallelism. In the book this
ratio was much closer to 1 (visually CPU was at about 85\% of
GPU-alone, i.e.\ \emph{faster} than the GPU); the explanation is
straightforward---a Ryzen 3900X has twice as many threads as an
i7-9750H, while the two GPUs (1660 Ti vs.\ 2060 Super) are in the same
Turing generation with similar fp64 ratios. Doubling the CPU's width
roughly doubles its throughput on this embarrassingly-parallel workload,
flipping the comparison.

\subsection{The Hybrid Curve Has the Same Shape, Just Flatter}

Both the book's curve and the replicated curve descend monotonically as
CPU workers are added to the GPU baseline, and both flatten as the
returns from each additional worker shrink. The marginal contribution
of CPU workers, however, differs substantially:

\begin{table}[H]
    \centering
    \caption{Marginal contribution of each additional CPU worker added
             to the GPU thread (this work).}
    \label{tab:marginal}
    \begin{tabular}{lrrr}
        \hline
        \textbf{Step} & \textbf{Time (s)} & \textbf{$\Delta$ time (s)} & \textbf{$\Delta$\,\% of GPU-alone} \\
        \hline
        GPU            & 76.54 &   ---  &   --- \\
        \,+1 CPU       & 68.45 & -8.09  & -10.6\% \\
        \,+2 CPU       & 64.40 & -4.05  &  -5.3\% \\
        \,+4 CPU       & 60.43 & -3.97  &  -5.2\% \\
        \,+8 CPU       & 56.19 & -4.24  &  -5.5\% \\
        \,+11 CPU      & 53.62 & -2.58  &  -3.4\% \\
        \hline
    \end{tabular}
\end{table}

The first CPU worker is worth a full 10\% of GPU-alone time. The next
ten together are worth another 20\%, and the last three workers (going
from 8~CPU to 11~CPU) buy only a 3.4\% improvement. This is the
\emph{Amdahl-like saturation of the hybrid}: once the work-queue
overhead per item and the fundamentally sequential portions of the
program (frame setup, image saves, the last-region tail of each frame)
become significant relative to the per-worker contribution, additional
workers cannot help. On the book's machine the same saturation appears
later (the curve is still descending visibly at +16~CPU and only flattens
at +23~CPU), which is consistent with having twice as many threads to
absorb.

\subsection{Why CPU-Only Is So Much Slower Here}

The CPU-only bar on this machine (163.5~s) is more than twice the
GPU-only bar (76.5~s). A back-of-envelope check using the per-region
statistics from the prior NVTX/CUDA-event profiling pass:

\begin{itemize}
    \item Total leaf regions, hybrid mode: 4{,}603
    \item Average CPU time per region: 535~ms (from
    \texttt{printCPUSummary} totals)
    \item Average GPU time per region: 13.5~ms
\end{itemize}

In CPU-only mode, all 4{,}603 leaves are routed to CPU workers. With 12
threads sharing the work and an average of 535~ms per region the
naively-predicted runtime is
$4603 \times 0.535 / 12 \approx 205\,\text{s}$. The measured 163.5~s
falls below this estimate because (a) CPU-only mode does not have the
$\sim$15~s outlier region described in the prior profiling report
disproportionately landing on one thread, since with 12 CPU threads
splitting all work the outlier no longer happens to be the tail
constraint; and (b) the corner-evaluation overhead in
\texttt{examine()}---which was $\sim$24~s across 6{,}104 calls in the
hybrid run---is amortized across more threads here. Either way, the
order-of-magnitude prediction is correct, confirming that the measured
makespan is consistent with the per-region measurements from the prior
pass.

\subsection{Speed-Up Summary}

Three speed-up numbers worth recording:

\[
\text{Speed-up of GPU over CPU12: } \frac{163.5}{76.5} = 2.14\times
\]
\[
\text{Speed-up of best hybrid over CPU12: } \frac{163.5}{53.6} = 3.05\times
\]
\[
\text{Speed-up of best hybrid over GPU-alone: } \frac{76.5}{53.6} = 1.43\times
\]

The hybrid is unambiguously the right choice on this platform; using
only the GPU leaves $\sim$30\% of performance on the table.

% -----------------------------------------------------------------------
\section{Comparison with the Book}

The qualitative finding from the book---``combining both
[CPU and GPU] cuts in half the execution time that either of them can
achieve''---does not translate verbatim to this machine, because on
this machine the CPU and GPU are not roughly on par. The corrected
statement for this platform is:

\begin{quote}
\emph{On a CPU/GPU pairing where the GPU is roughly 2$\times$ faster
than the CPU at full thread parallelism, combining both cuts the
GPU-alone execution time by another 30\% and the CPU-alone execution
time by a factor of 3.}
\end{quote}

The structural conclusion from Section~11.5.1 of the book---that
dynamic load balancing via the shared work queue is effective and
delivers monotone improvements as more workers are added---is fully
supported. The hybrid curves have the same shape, the same convexity,
and the same saturation behavior; only the absolute baseline and the
slope differ, and both differences are explained by the difference in
CPU width (12 vs.\ 24 threads).

It is worth noting one residual asymmetry. The book's curve appears to
flatten at +16~CPU, with +23~CPU offering little further improvement;
on a 24-thread machine this is consistent with two effects compounding:
the marginal contribution of each worker shrinking, and SMT
contention---hyper-threads do not double pure-compute throughput
because the integer/floating-point pipelines are shared per physical
core. On this machine, the same effect is present (the i7-9750H is
6c/12t, also 2 hyperthreads per core) but the configurations sampled
stop at +11~CPU, which is the last point before saturation rather than
deep into it.

% -----------------------------------------------------------------------
\section{Follow-Up Experiments}

The instrumentation patch left two unused affordances that are
straightforward to exploit:

\begin{enumerate}
    \item \textbf{\texttt{SAVE=0} sweep.} Re-running with
    \texttt{SAVE=0 ./sweep.sh} would isolate the pure-compute portion
    of the runtime. The PNG-encode tail is one of the sequential
    portions limiting hybrid speed-up; quantifying its share would
    reveal whether eliminating it (e.g.\ moving image encoding into a
    background writer thread) is worth the engineering cost.
    \item \textbf{\texttt{diffT} sensitivity.} The prior profiling
    report identified a 15.3-second worst-case CPU region that was
    misclassified as uniform by the corner heuristic
    (\texttt{diffT}\,=\,0.5). \texttt{DIFFT=0.2 ./sweep.sh} would test
    whether tightening the threshold improves either the worst-case CPU
    region's wall time or the overall hybrid total. The trade-off is
    more queue traffic from forced splits.
\end{enumerate}

Two further parameters not exposed by the harness but worth varying
in a future pass:

\begin{itemize}
    \item \textbf{CUDA block size.} \texttt{BLOCK\_SIDE} is hard-coded
    to 16 in \texttt{kernel.cu} (a 256-thread 2D block). Trying 8 and
    32 would test occupancy-vs-divergence tradeoffs.
    \item \textbf{Work-queue policy.} The existing
    \texttt{MandelRegion::Compare} functor already encodes a
    largest-first ordering; swapping the \texttt{deque} for a
    \texttt{priority\_queue} would deliver large regions to the GPU
    thread first, which is where they are processed fastest, and
    might recover some of the lost speed-up on the right-hand
    saturation tail.
\end{itemize}

% -----------------------------------------------------------------------
\section{Conclusion}

The hybrid Mandelbrot renderer from Barlas (2022) reproduces, in
spirit, on an i7-9750H + GTX~1660~Ti~Max-Q laptop: the curve has the
same monotonic shape, the dynamic work queue delivers good load
balance, and the best hybrid is the best configuration on the machine.
The book's specific quantitative claim that hybrid execution
``cuts in half'' the execution time of either CPU or GPU alone is a
property of a CPU/GPU pairing where the two are roughly equal in
throughput; on a system where the GPU is the dominant accelerator, the
hybrid still wins but by a smaller margin (1.43$\times$ over GPU-alone
rather than $\sim$2$\times$).

The instrumentation patch needed to extract the metric Fig~11.15
actually plots was small (an \texttt{extern "C"} flag pair, a
\texttt{clock\_gettime} bracket, three printf gates) and orthogonal to
the existing NVTX/CUDA-event work. The batch driver and plot script are
generic enough to re-use for any future thread-configuration or
threshold sweep, simply by changing the \texttt{CONFIGS} array or the
environment variables.

\end{document}
